% !TEX root = ../PRACA_MAIN.tex

\section{Wprowadzenie do problematyki systemów rekomendacyjnych}

Rozdział pierwszy przedstawia teoretyczne oraz rynkowe podstawy systemów rekomendacyjnych. Omówiono w nim zjawisko przeciążenia poznawczego na rynku platform VOD, formalny podział metod personalizacji treści oraz technologiczną ewolucję algorytmów -- od klasycznych podejść opartych na algebrze macierzy po zaawansowane struktury grafowe. Rozdział kończy analiza kluczowych wyzwań wdrożeniowych związanych z projektowaniem tego typu systemów.

\subsection{Zjawisko przeciążenia poznawczego i specyfika rynku VOD} % REVIEWED

Rozwój technologii internetowych w pierwszych dekadach XXI wieku doprowadził do cyfryzacji głównych segmentów rynku rozrywkowego. W sektorze wideo obserwuje się stopniowe odchodzenie od tradycyjnej telewizji na rzecz platform VOD. Rynek ten zdominowały serwisy globalne (np. Netflix, Max, Amazon Prime Video) oraz dostawcy lokalni (np. Player.pl, TVP VOD). Analogiczne zmiany nastąpiły w branży muzycznej, gdzie czołową pozycję zajmują platformy takie jak Spotify i Apple Music. Model subskrypcyjny wprowadzono także na rynku wydawniczym, umożliwiając stały dostęp do e-booków oraz audiobooków. W rezultacie użytkownicy mogą korzystać z obszernych baz multimedialnych z poziomu pojedynczej aplikacji.

Rosnąca objętość udostępnianych katalogów stanowi jednak istotny problem z perspektywy użytkownika. Ze względu na znaczną liczbę dostępnych materiałów, odbiorcy często doświadczają zjawiska przeciążenia poznawczego (ang. \foreignlanguage{english}{\textit{cognitive overload}}) \footcite{bollenUnderstandingChoiceOverload2010}. Zjawisko to występuje powszechnie w całym sektorze cyfrowej rozrywki, co potwierdzają dane rynkowe. Baza IMDb indeksuje ponad 30 milionów tytułów filmowych i telewizyjnych \footcite{PressRoomIMDb}. Katalog platformy Spotify obejmuje ponad 100 milionów utworów \footcite{AboutSpotify}, natomiast serwis BookBeat oferuje subskrybentom dostęp do ponad miliona pozycji literackich \footcite{BookBeatONas}.

Setki tysięcy produkcji w ramach pojedynczego serwisu wywołują u odbiorcy efekt tzw. paradoksu wyboru \footcite{schwartzParadoxChoiceWhy2005}. Nadmierna ekspozycja na treści przytłacza użytkownika, sprawiając, że całościowa eksploracja katalogu staje się fizycznie niemożliwa, a podjęcie decyzji o wyborze seansu wiąże się ze znacznym wysiłkiem poznawczym. Takie przeciążenie w wielu przypadkach prowadzi do zniechęcenia, frustracji widza, a ostatecznie do rezygnacji z subskrypcji. Rozwiązania tradycyjne, takie jak wyszukiwanie frazowe czy statyczne, redakcyjne zestawienia popularności, okazują się niewystarczające, ponieważ wymagają one aktywnej inicjatywy ze strony odbiorcy lub promują pozycje najpopularniejsze, nieuwzględniające jego indywidualnych preferencji \footcite{krawczukSystemyRekomendacyjne2024}.

W konsekwencji zaawansowane systemy rekomendacyjne przestały być jedynie opcjonalnym dodatkiem funkcjonalnym, stając się integralnym i kluczowym elementem architektury współczesnych platform cyfrowych \footcite[23--24]{Bernardelli2024}. Ich nadrzędnym zadaniem jest ciągłe, zautomatyzowane analizowanie ogromnych katalogów treści i filtrowanie ich w taki sposób, aby skutecznie chronić widza przed zjawiskiem paraliżu decyzyjnego. Zamiast zmuszać odbiorcę do samodzielnego przeszukiwania tysięcy tytułów, algorytmy te dostarczają spersonalizowane propozycje, ściśle dopasowane do jego profilu, historii oglądania oraz preferencji. W cyfrowym modelu biznesowym, opartym na miesięcznych subskrypcjach, to właśnie trafność tych podpowiedzi stanowi podstawowe narzędzie budujące lojalność i długoterminowe zaangażowanie użytkowników. Skuteczność takiego podejścia najlepiej obrazują wewnętrzne dane platformy Netflix, z których wynika, że precyzyjne sugestie algorytmiczne odpowiadają za wygenerowanie aż 80\% całkowitego czasu, jaki widzowie spędzają na odtwarzaniu materiałów w serwisie \footcite{gomez-uribeNetflixRecommenderSystem2015}.

\subsection{Definicje, formalizm matematyczny i struktura danych w personalizacji treści} % REVIEWED

W niniejszym podrozdziale wprowadzono podstawowe pojęcia pozwalające zrozumieć działanie systemów personalizacji treści. Przedstawiono rozwój oraz matematyczny zapis problemu rekomendacji, a także omówiono specyfikę danych wykorzystywanych przez współczesne algorytmy uczenia maszynowego.

\subsubsection{Ewolucja pojęcia systemów rekomendacji}

Historycznie pojęcie systemów rekomendacji ewoluowało wraz z postępującą cyfryzacją oraz narastającym problemem przeciążenia poznawczego. Pierwsze rozwiązania badawcze miały charakter ściśle kooperacyjny i opierały się na świadomej współpracy odbiorców. Pionierskim przykładem jest stworzony w 1992 roku przez badaczy z Xerox Palo Alto Research Center (PARC) system Tapestry \footcite{goldbergUsingCollaborativeFiltering1992}. Powstał on w odpowiedzi na dynamicznie rosnącą liczbę wiadomości w sieciach korporacyjnych i na grupach dyskusyjnych, proponując nowatorskie podejście do selekcji strumieni danych. Tapestry umożliwiał użytkownikom filtrowanie przychodzących dokumentów na podstawie jawnych, manualnych adnotacji i ocen wprowadzonych wcześniej przez innych czytelników (np. żądanie wyświetlania wyłącznie tych wiadomości, które konkretny współpracownik oznaczył jako wartościowe). To właśnie w publikacji opisującej ten system po raz pierwszy w literaturze sformułowano termin "filtrowanie kolaboracyjne", kładąc podwaliny pod cały współczesny nurt badań nad personalizacją treści. Chociaż system Tapestry działał zadowalająco, to jednak ze względu na wymóg budowania skomplikowanych zapytań i ograniczenie do małych, zamkniętych społeczności, w których użytkownicy znali się nawzajem, jego koncepcja szybko ewoluowała w stronę w pełni zautomatyzowanych, bezobsługowych algorytmów na dużą skalę.

Przełomowym momentem dla rozwoju pojęć w tej dziedzinie była praca opublikowana w 1997 roku przez P. Resnicka i H. Variana \footcite{RecommenderSystems1997}. Badacze ci zaproponowali, aby wąskie pojęcie "filtrowania kolaboracyjnego" zastąpić szerszym i bardziej uniwersalnym terminem "systemy rekomendacyjne". Argumentowali oni, że nowa generacja algorytmów nie wymaga już celowej współpracy między użytkownikami, którzy w środowisku sieciowym mogą pozostawać dla siebie całkowicie anonimowi. Zgodnie z ich ujęciem, system rekomendacyjny to dowolne zautomatyzowane narzędzie, które agreguje oceny -- zarówno te wprowadzane jawnie, jak i te pozyskane w sposób niejawny z zachowania odbiorców (ang. \foreignlanguage{english}{\textit{implicit feedback}}), a następnie kieruje trafne sugestie do odpowiednich osób.

Wraz z dynamicznym rozwojem metod uczenia maszynowego na początku XXI wieku, definicja ta uległa dalszemu rozszerzeniu i abstrakcji. Klasyczne już dziś ujęcie zaproponował w 2002 roku R. Burke, definiując system rekomendacji jako każdy program, którego nadrzędnym zadaniem jest generowanie zindywidualizowanych sugestii lub aktywne kierowanie użytkownika do interesujących go obiektów w obrębie bardzo rozległej przestrzeni opcji \footcite{burkeHybridRecommenderSystems2002}. To pojemne ujęcie usankcjonowało rozwój kolejnych architektur, w tym filtrowania opartego na treści (ang. \foreignlanguage{english}{\textit{Content-Based Filtering}} -- CBF) oraz wysoce złożonych modeli hybrydowych. Współcześnie, za sprawą algorytmów głębokiego uczenia, zadanie to uległo radykalnej automatyzacji. W ekosystemach cyfrowych platform VOD systemy rekomendacyjne przestały być jedynie narzędziem pomocniczym wspomagającym wybór, stając się ich centralnym elementem i kształtując codzienne decyzje użytkowników w znacznie większym stopniu niż tradycyjne wyszukiwanie czy statyczne rankingi popularności.

\subsubsection{Formalizacja matematyczna problemu rekomendacji}

Systemy rekomendacyjne są nierozerwalnie związane z problemem badawczym, który można sformalizować jako problem optymalizacyjny z wykorzystaniem funkcji użyteczności \footcite[36--37]{malinowskiAlgorytmRekomendacjiBazujacy2022}:
\begin{equation}
    u \colon C \times O \to W,
\end{equation}
gdzie:
\begin{itemize}
    \item $C = \{c_1, c_2, \dots, c_m\}$ -- skończony zbiór użytkowników wykorzystujących system rekomendacyjny,
    \item $O = \{o_1, o_2, \dots, o_n\}$ -- skończony zbiór wszystkich możliwych obiektów podlegających rekomendacji,
    \item $u$ -- funkcja użyteczności mierząca przydatność obiektu dla użytkownika,
    \item $W$ -- uporządkowany zbiór wartości (np. zbiór nieujemnych liczb całkowitych).
\end{itemize}

Głównym zadaniem algorytmów rekomendacyjnych jest znalezienie obiektów o najwyższej użyteczności dla danego odbiorcy. Matematycznie sprowadza się to do wyznaczenia dla każdego użytkownika optymalnego zbioru rekomendacji \footcite{adomaviciusNextGenerationRecommender2005a}:
\begin{equation}
    R_c^* = \operatorname*{arg\,max}_{o \in O} u(c,o),
\end{equation}
gdzie:
\begin{itemize}
    \item $R_c^*$ -- zbiór elementów z przestrzeni $O$ maksymalizujących funkcję użyteczności,
    \item $u(c,o)$ -- użyteczność obiektu $o$ dla konkretnego użytkownika $c$.
\end{itemize}

Jednym z głównych wyzwań związanych z działaniem systemów personalizacji treści jest zjawisko rzadkości danych (\foreignlanguage{english}{\textit{data sparsity}}), wynikające z dysponowania nielicznymi informacjami o użytkownikach lub z ich całkowitego braku \footcite{mcauleyHiddenFactorsHiddenPazdziernik122013}. Taka sytuacja ma miejsce w przypadku, gdy użytkownik po raz pierwszy ma kontakt z systemem (tzw. problem zimnego startu) lub gdy system celowo nie gromadzi danych historycznych w związku z wymogami ochrony prywatności. Z tego powodu klasyczna postać problemu rekomendacji bywa często modyfikowana w taki sposób, aby algorytm opierał się bezpośrednio na podobieństwie obiektów, a nie na złożonym profilu użytkownika. Wymaga to zastąpienia funkcji użyteczności odwzorowaniem $w$ \footcite{malinowskiZastosowanieGrafowSieci2021}:
\begin{equation}
    w \colon O \times O \to W,
\end{equation}
gdzie:
\begin{itemize}
    \item $w$ -- nowa funkcja określająca siłę powiązania między dwoma obiektami.
\end{itemize}

Dzięki temu pierwotny problem wyznaczenia rekomendacji dla użytkownika transformuje się w optymalizację zbioru dla konkretnego elementu. W konsekwencji zbiór polecanych przedmiotów definiuje się następująco \footcite{malinowskiZastosowanieGrafowSieci2021}:
\begin{equation}
    R_{o^*}^* = \operatorname*{arg\,max}_{o \in O} w(o^*, o),
\end{equation}
gdzie:
\begin{itemize}
    \item $R_{o^*}^*$ -- optymalny zbiór rekomendowanych przedmiotów powiązanych z obiektem startowym,
    \item $o^*$ -- ustalony przez użytkownika obiekt bazowy (np. film, w którego miniaturę właśnie kliknął),
    \item $w(o^*, o)$ -- siła powiązania (matematyczne podobieństwo) między ocenianym obiektem $o \in O$ a obiektem bazowym $o^*$.
\end{itemize}

\subsubsection{Paradygmaty modelowania: przewidywanie ocen, a prognoza rankingu} % REVIEWED

Wyniki działania algorytmów rekomendacyjnych prezentowane są użytkownikowi na podstawie dwóch głównych paradygmatów \footcite{malinowskiZastosowanieGrafowSieci2021}:
\begin{enumerate}
    \item Przewidywanie ocen: model prognozuje konkretną wartość liczbową oceny (np. 4,5/5 gwiazdek). Celem jest tutaj minimalizacja błędu bezwzględnego estymacji (np. RMSE) między przewidywaną oceną a rzeczywistą. Taki paradygmat stanowił fundament historycznego konkursu \foreignlanguage{english}{\textit{Netflix Prize}} (2006--2009) i algorytmu \foreignlanguage{english}{\textit{Cinematch}}.
    \item Prognoza rankingu: system ignoruje absolutną wartość oceny na rzecz wygenerowania uporządkowanej listy $N$ najlepszych pozycji, dopasowanych specyficznie do użytkownika. Optymalizacji podlega względna kolejność proponowanych elementów. Podejście to zostało pierwotnie spopularyzowane przez wyszukiwarki internetowe oraz pionierskie platformy e-commerce (np. Amazon), dla których znacznie cenniejsze od predykcji bezwzględnej oceny produktu było zaprezentowanie klientowi zoptymalizowanej listy najbardziej trafnych propozycji zakupowych. Z czasem algorytmy te stały się nieodzownym elementem platform internetowych, w tym serwisów streamingowych \footcite{deshpandeItembasedTopNRecommendationStyczen012004}.
\end{enumerate}
W nowoczesnych systemach rozrywkowych VOD zdecydowanie preferuje się prognozy rankingowe (Top-N). Jak wykazuje Harald Steck \footcite{steckEvaluationRecommendationsRatingprediction2013}, z perspektywy biznesowej oraz doświadczeń użytkownika, kluczowe jest zaprezentowanie najtrafniejszych obiektów na samym szczycie listy, a nie przewidywanie dokładnych ocen liczbowych. Odejście od klasycznej predykcji ocen wynika również z masowego przejścia platform na gromadzenie danych niejawnych. Współcześni odbiorcy rzadko udzielają bezpośrednich ocen punktowych, co wymusza na algorytmach analizę behawioralną (np. mierzenie czasu odtwarzania materiału). Dane tego typu nie przekładają się wprost na bezwzględną ocenę liczbową, jednak idealnie wpisują się w zadanie optymalizacji względnego rankingu proponowanych treści.

\subsubsection{Struktura danych i źródła informacji zwrotnej} % REVIEWED

Systemy rekomendacyjne to w swej istocie systemy przetwarzania informacji, których skuteczność i architektura zależą bezpośrednio od wolumenu, jakości oraz wielowymiarowości gromadzonych danych wejściowych. Dane te mogą mieć różną formę i pochodzenie, a stopień ich wykorzystania determinuje wybór odpowiedniej techniki rekomendacji. W literaturze przedmiotu strukturę danych zasilających systemy rekomendacyjne dzieli się najczęściej na trzy główne obiekty: przedmioty, użytkowników oraz transakcje \footcite{ricciRecommenderSystemsHandbook2010}:

\begin{itemize}
    \item Przedmioty: obiekty podlegające rekomendacji charakteryzują się różnym stopniem złożoności oraz określoną wartością (użytecznością) dla odbiorcy, z którą wiąże się koszt poznawczy lub finansowy. W kontekście portali streamingowych (VOD) pozycje mogą być reprezentowane w sposób minimalistyczny (wyłącznie jako unikalny identyfikator) lub za pomocą bogatych zbiorów metadanych. Obejmują one atrybuty katalogowe (np. gatunki, reżyserzy, obsada), a także zaawansowane cechy generowane automatycznie przez systemy analizy obrazu i dźwięku (np. VCR -- ang. \foreignlanguage{english}{\textit{Video Content Recognition}}), takie jak tempo montażu czy tonacja kolorystyczna.

    \item Użytkownicy: model użytkownika gromadzi informacje pozwalające na personalizację rekomendacji, kodując jego preferencje i potrzeby. W zależności od wybranego algorytmu reprezentacja ta może przyjmować różną formę. W systemach demograficznych profil opiera się na statycznych parametrach (np. wiek, płeć), podczas gdy w filtrowaniu kolaboracyjnym użytkownik modelowany jest najczęściej poprzez wektor ukrytych cech lub bezpośrednio przez historię ocen. Profil uzupełniany jest często o kontekst dynamiczny (np. używane urządzenie, pora dnia, geolokalizacja) oraz dane behawioralne odzwierciedlające wzorce nawigacji w systemie.

    \item Transakcje: zapisy interakcji pomiędzy użytkownikiem a systemem rekomendacyjnym. Stanowią one sekwencyjne logi zachowań, które obok samego faktu wyboru przedmiotu mogą zawierać informacje o kontekście zapytania oraz ewentualnej informacji zwrotnej. To właśnie na podstawie danych transakcyjnych algorytmy uczą się wzorców preferencji.
\end{itemize}

Z punktu widzenia optymalizacji matematycznej modelu i budowy funkcji celu, kluczowy jest podział informacji zwrotnej (reprezentującej historyczne transakcje) na dwie główne kategorie:

\begin{itemize}
    \item Jawne oceny (ang. \foreignlanguage{english}{\textit{explicit feedback}}): bezpośrednie wyrażenie preferencji przez użytkownika. Przyjmują one najczęściej formę ocen numerycznych (np. skala 1--5 gwiazdek), ocen porządkowych, czy też wartości binarnych (np. pozytywny/negatywny). Choć dostarczają precyzyjnych i jednoznacznych informacji o upodobaniach, ich główną wadą jest wysoki koszt poznawczy wymagany od użytkownika, co sprawia, że występują w systemie skrajnie rzadko, prowadząc do silnego problemu rzadkości macierzy danych.

    \item Dorozumiane ślady aktywności (ang. \foreignlanguage{english}{\textit{implicit feedback}}): zachowania wywnioskowane w sposób pośredni z akcji użytkownika (np. czas oglądania, historia odtworzeń, przewijanie wideo, liczba kliknięć). Są to dane łatwe do masowego gromadzenia, jednak charakteryzują się sporym zaszumieniem. Modelowanie informacji dorozumianych wiąże się z problematyką jednoklasowego filtrowania kolaboracyjnego \footcite{rendleBPRBayesianPersonalized2009}. Główne wyzwanie polega na interpretacji braku interakcji: w przeciwieństwie do ocen jawnych, brak interakcji z filmem nie musi oznaczać niechęci do niego, lecz bardzo często wynika z faktu, że użytkownik nie był świadomy jego istnienia.
\end{itemize}

Współczesne środowiska wielkoskalowych systemów VOD, z uwagi na marginalny odsetek ocen jawnych, opierają się w głównej mierze na wykorzystaniu dorozumianych śladów aktywności. Zdolność algorytmów do skutecznego wyciągania użytecznych informacji z tych obfitych, lecz przepełnionych zbędnymi informacjami zbiorów danych, przy jednoczesnym uwzględnieniu dostępnych metadanych, jest kluczowa dla jakości działania całego systemu.

\subsection{Ewolucja technologiczna paradygmatów rekomendacyjnych} % REVIEWED

Systemy rekomendacyjne funkcjonują i są rozwijane na rynku od ponad trzydziestu lat. Na przestrzeni tego czasu środowiska informatyczne przeszły całkowitą transformację, z czym wiązała się również dynamiczna zmiana preferencji, przyzwyczajeń i oczekiwań samych użytkowników. Te głębokie przemiany technologiczne i behawioralne wymusiły ciągłą ewolucję w sposobie projektowania algorytmów. Opierając się na współczesnym przeglądzie literatury przeprowadzonym przez Razę i in. \footcite{razaComprehensiveReviewRecommender2024}, historyczny rozwój paradygmatów rekomendacyjnych można podzielić na trzy główne ery. Odzwierciedlają one rosnącą złożoność w modelowaniu preferencji użytkowników oraz postęp w dziedzinie sztucznej inteligencji, przy czym każda z er charakteryzuje się odmiennym podejściem do analizy danych i reprezentacji wiedzy:

\vspace{0.5cm}
\begin{figure}[htbp]
    \centering
    \resizebox{\textwidth}{!}{
        \begin{tikzpicture}[
                level distance=2.8cm,
                level 1/.style={sibling distance=7.5cm},
                level 2/.style={sibling distance=3.5cm},
                level 3/.style={sibling distance=3.2cm},
                box/.style={draw=black!70, rectangle, rounded corners=3pt, align=center, inner sep=6pt},
                root/.style={box, fill=blue!15, font=\bfseries\Large},
                main/.style={box, fill=blue!5, font=\bfseries},
                leaf/.style={box, fill=white, font=\small, text width=3.0cm},
                edge from parent/.style={draw=black!70, thick, ->, >=stealth},
                edge from parent path={(\tikzparentnode.south) -- +(0,-0.6) -| (\tikzchildnode.north)}
            ]

            \node[root] {Systemy Rekomendacyjne}
            child {node[main] {Podejścia Tradycyjne}
                    child {node[leaf] {Filtrowanie Treściowe\\(Content-Based)}}
                    child {node[leaf] {Filtrowanie\\Kolaboracyjne (CF)}
                            child {node[leaf, fill=gray!5, font=\footnotesize, text width=2.6cm] {Oparte na pamięci\\($k$-NN)}}
                            child {node[leaf, fill=gray!5, font=\footnotesize, text width=2.6cm] {Oparte na modelu\\(SVD, BPR)}}
                        }
                }
            child {node[main] {Uczenie Głębokie\\(Deep Learning)}
                    child {node[leaf] {Modele Perceptronowe\\(np. NCF, NeuMF)}}
                    child {node[leaf] {Modele Sekwencyjne\\(np. SASRec)}}
                }
            child {node[main] {Podejścia Grafowe\\(GNN)}
                    child {node[leaf] {Grafy Interakcji\\(np. LightGCN)}}
                    child {node[leaf] {Grafy Wiedzy\\(np. KGCN)}}
                };
        \end{tikzpicture}
    }
    \caption{Taksonomia i historyczny podział algorytmów rekomendacyjnych uwzględnionych w badaniach.}
    \label{fig:taksonomia_rekomendacji}

    \vspace{0.2cm}
    {\raggedright \small \textbf{Źródło:} opracowanie własne.\par}
\end{figure}

\vspace{0.5cm}

\subsubsection{Tradycyjne systemy rekomendacyjne}

Historyczny rozwój systemów rekomendacyjnych ukształtowały trzy klasyczne paradygmaty, które do dziś odgrywają rolę standardowych modeli referencyjnych w badaniach naukowych \footcite{suSurveyCollaborativeFiltering2009}:
\begin{itemize}
    \item \textbf{Filtrowanie oparte na treści}: rekomenduje obiekty na podstawie ich podobieństwa do tych wybieranych w przeszłości (często obliczając iloczyn skalarny wektorów). Metoda ta jest stosunkowo odporna na problem „zimnego startu” w przypadku wprowadzania do katalogu nowych premier filmowych. Wynika to z faktu, że algorytm do działania nie wymaga zgromadzenia początkowej masy interakcji od społeczności. Wystarczy sama obecność metadanych opisujących nowy film (np. gatunek, obsada, reżyser), by móc od razu dopasować go do profili widzów. Z drugiej strony, ścisłe uzależnienie predykcji od dotychczasowej historii oglądania prowadzi do zjawiska nadmiernej specjalizacji. System nie posiada mechanizmu eksploracji i promuje wyłącznie pozycje podobne do tych już obejrzanych. W efekcie zamyka to użytkownika w bańce filtracyjnej, pozbawiając go możliwości odkrywania nowych gatunków i tytułów. Produkcje te, choć mogłyby okazać się dla niego atrakcyjne, przez swoje niewystarczające podobieństwo do dotychczasowej historii oglądania są całkowicie pomijane przez algorytm.
    \item \textbf{Filtrowanie kolaboracyjne:} opiera się na analizie opinii i zachowań całej społeczności. Główne założenie tego podejścia mówi, że użytkownicy o podobnej historii interakcji (np. oglądający te same filmy) będą mieli zbliżone preferencje w przyszłości. Klasycznie algorytmy te dzieli się na metody oparte na pamięci wykorzystujące m.in. algorytm $k$-najbliższych sąsiadów oraz znacznie wydajniejsze metody oparte na modelach. Podejścia pamięciowe, choć intuicyjne, wykazują poważne problemy ze skalowalnością i drastycznie tracą na skuteczności przy bardzo rzadkich danych. Odpowiedzią na te ograniczenia stała się technika faktoryzacji macierzy (m.in. rozkład według wartości osobliwych, ang. \foreignlanguage{english}{\textit{Singular Value Decomposition}} -- SVD, oraz ALS), która zrewolucjonizowała uczenie modeli bazowych. Pozwala ona na rzutowanie rzadkiej macierzy interakcji do gęstej przestrzeni tzw. czynników ukrytych (\foreignlanguage{english}{\textit{latent factors}}), dzięki czemu system potrafi uczyć się nieoczywistych wzorców. Należy jednak zaznaczyć, że w przeciwieństwie do metod opartych na treści, czyste filtrowanie kolaboracyjne jest całkowicie podatne na problem „zimnego startu” -- nie potrafi trafnie polecić nowo dodanego filmu, dopóki ten nie zbierze krytycznej masy odtworzeń od początkowych widzów.
    \item \textbf{Modele hybrydowe:} stanowią pragmatyczny kompromis łączący podejścia kolaboracyjne i treściowe, zaprojektowany z myślą o niwelowaniu słabości pojedynczych paradygmatów. Poprzez połączenie obu technik system potrafi jednocześnie przełamać bańkę filtracyjną (wykorzystując eksploracyjny potencjał CF) oraz skutecznie rekomendować nowości kinowe nieposiadające historii odtworzeń (wykorzystując metadane z CBF). Klasyczną taksonomię takich rozwiązań przedstawił R. Burke \footcite{burkeHybridRecommenderSystems2002}, wyróżniając kluczowe strategie ich integracji. Należą do nich m.in.: łączenie wyników poprzez ważone sumowanie predykcji niezależnych modeli, wykorzystywanie wyników jednego algorytmu jako atrybutów dla drugiego, wieloetapowe kaskadowanie predykcji oraz dynamiczne przełączanie (\foreignlanguage{english}{\textit{switching}}). Ta ostatnia strategia pozwala systemowi np. zastosować model oparty na treści dla najnowszych premier, a płynnie przejść na model kolaboracyjny w przypadku filmów posiadających już bogatą, historyczną reprezentację w bazie danych.
\end{itemize}

\subsubsection{Era Uczenia Głębokiego}

Adaptacja głębokich sieci neuronowych w systemach rekomendacyjnych pozwoliła na rozwiązanie części ograniczeń metod klasycznych. Nowe modele wprowadziły przede wszystkim mechanizmy automatycznego uczenia się reprezentacji (ang. \foreignlanguage{english}{\textit{representation learning}}) oraz modelowania nieliniowych zależności między użytkownikiem a przedmiotem \footcite{zhangDeepLearningBased2020}. Rozwój tej dziedziny opierał się na zastosowaniu kilku głównych architektur:

\begin{itemize}
    \item \textbf{Wielowarstwowe perceptrony (ang. \foreignlanguage{english}{\textit{Multi-Layer Perceptron}} -- MLP):} uogólniły metody oparte na faktoryzacji macierzy. W modelu \foreignlanguage{english}{\textit{Neural Collaborative Filtering}} (NCF) klasyczny iloczyn skalarny zastąpiono siecią neuronową, co pozwoliło na aproksymację funkcji nieliniowych i zwiększyło pojemność predykcyjną algorytmu.

    \item \textbf{Autoenkodery i sieci splotowe (VAE / CNN):} wariacyjne autoenkodery (ang. \foreignlanguage{english}{\textit{Variational Autoencoders}} -- VAE) znalazły zastosowanie w redukcji wymiarowości przy dużej rzadkości danych. Sieci splotowe (ang. \foreignlanguage{english}{\textit{Convolutional Neural Networks}} -- CNN) posłużyły natomiast do analizy multimediów (np. miniatur lub klatek wideo), pozwalając na bezpośrednie włączenie cech wizualnych do modelu.

    \item \textbf{Modele sekwencyjne i Transformery:} do analizy chronologicznej historii odtworzeń początkowo wykorzystywano sieci rekurencyjne (ang. \foreignlanguage{english}{\textit{Recurrent Neural Networks}} -- RNN). Z czasem zastąpiono je mechanizmami atencji i architekturą Transformer (np. modele SASRec czy BERT4Rec), które lepiej radzą sobie z modelowaniem zarówno krótko-, jak i długoterminowych preferencji użytkownika.
\end{itemize}

\subsubsection{Zaawansowane modele systemów rekomendacyjnych}

Współczesne modele rekomendacyjne odchodzą od płaskich, tabelarycznych reprezentacji użytkownika na rzecz struktur opartych na grafach i relacjach \footcite{malinowskiZastosowanieGrafowSieci2021}:
\begin{itemize}
    \item \textbf{Systemy oparte na sesjach (ang. \foreignlanguage{english}{\textit{Session-Based Recommendation}} -- SBR):} modelują krótkoterminowe preferencje użytkownika w ramach pojedynczej, często anonimowej sesji, wykorzystując łańcuchy Markowa lub analizę sekwencyjną.

    \item \textbf{Grafowe sieci neuronowe (ang. \foreignlanguage{english}{\textit{Graph Neural Networks}} -- GNN):} reprezentują system VOD jako graf dwudzielny lub heterogeniczny. Modele oparte na propagacji osadzeń (takie jak architektura LightGCN) pozwalają na wydajne modelowanie relacji wyższego rzędu pomiędzy użytkownikami a przedmiotami. Z kolei sieci splotowe na grafach wiedzy (ang. \foreignlanguage{english}{\textit{Knowledge Graph Convolutional Networks}} -- KGCN) uwzględniają bogaty kontekst semantyczny. Wspiera to koncepcję wyjaśnialnej sztucznej inteligencji (ang. \foreignlanguage{english}{\textit{Explainable AI}} -- XAI), pozwalając generować komunikaty typu: „Polecamy ten film, ponieważ polubiłeś reżysera X, który pracował także przy produkcji Y”.

    \item \textbf{Uczenie ze wzmocnieniem (ang. \foreignlanguage{english}{\textit{Reinforcement Learning}} -- RL):} traktuje problem rekomendacji jako sekwencyjny proces decyzyjny agenta. Algorytm optymalizuje kompromis między eksploatacją znanych treści a eksploracją nowości, maksymalizując długoterminową satysfakcję użytkownika, zamiast skupiać się wyłącznie na natychmiastowej nagrodzie (np. pojedynczym kliknięciu).
\end{itemize}
Przedstawiona taksonomia obrazuje ewolucję algorytmów rekomendacyjnych: od prostych heurystyk, przez gęste modele nieliniowe, aż po paradygmaty oparte na grafach. Wybór reprezentatywnych algorytmów z różnych gałęzi tej klasyfikacji pozwoli w dalszej części pracy na empiryczne porównanie skuteczności modeli pochodzących z różnych generacji. Zanim jednak zdefiniowana zostanie metodyka eksperymentów, niezbędne jest przyjrzenie się najnowszemu dorobkowi naukowemu w badanej dziedzinie.

\subsection{Przegląd literatury naukowej} % REVIEWED

W niniejszym podrozdziale dokonano przeglądu literatury naukowej z zakresu systemów rekomendacyjnych, ze szczególnym uwzględnieniem podejść opartych na filtrowaniu kolaboracyjnym, sztucznych sieciach neuronowych oraz grafowych strukturach wiedzy. Analiza obejmuje zarówno klasyczne metody faktoryzacji, jak i najnowsze osiągnięcia z obszaru głębokiego uczenia w odniesieniu do specyficznych uwarunkowań rynku VOD.

Przegląd literatury naukowej w kontekście usług streamingowych warto rozpocząć od analizy konkursu Netflix Prize (2006--2009) omówionego przez Hallinana i Striphasa \footcite{hallinanRecommendedYouNetflix2016}. Autorzy wykazują, że publiczne udostępnienie zbioru obejmującego ponad 100 milionów ocen przypisanych przez 480 tysięcy użytkowników dla 17 tysięcy filmów stanowiło punkt zwrotny w rozwoju algorytmów rekomendacyjnych. Rywalizacja o poprawę predykcji błędu RMSE o 10\% względem algorytmu Cinematch udowodniła, że tradycyjne metody oparte na sąsiedztwie są niewystarczające w środowisku o takiej skali, głównie ze względu na wysoki narzut pamięciowy i brak skalowalności obliczeniowej. To właśnie konkurs Netflixa spopularyzował modele czynników ukrytych, w tym rozkład SVD. Hallinan i Striphas zwracają również uwagę na skomplikowane zjawiska polaryzacyjne, określane jako wyzwanie \textit{Napoleon Dynamite problem}, w którym przewidywanie gustów użytkowników dla niszowych treści wymagało odrzucenia pojedynczych algorytmów na rzecz złożonych zespołów modeli hybrydowych.

Szczegółowe omówienie podstaw teoretycznych oraz wyzwań wdrożeniowych przedstawia praca Krawczuka \footcite{krawczukSystemyRekomendacyjne2024}. Autor przeprowadza analizę poszczególnych paradygmatów rekomendacji, oceniając ich przydatność w warunkach produkcyjnych. W pierwszej kolejności analizuje techniki CBF, wykazując, że ich podstawową zaletą jest niezależność od aktywności reszty społeczności oraz odporność na rzadkość ocen innych użytkowników. Zauważa jednak ograniczenia natury poznawczej: algorytmy CBF wykazują tendencję do nadmiernej specjalizacji, zamykając użytkownika w tzw. bańce filtracyjnej i utrudniając odkrywanie nieoczywistych, lecz trafnych pozycji, a także zmagają się z problemem zimnego startu dla nowych obiektów bez uzupełnionych metadanych. Przechodząc do podejść typu CF, Krawczuk akcentuje, że wykorzystywanie zbiorowych zachowań użytkowników pozwala na rekomendowanie treści niezależnie od ich opisu semantycznego. Metody te charakteryzują się jednak wysoką wrażliwością na rzadkość ocen, gdzie wskaźnik wypełnienia macierzy interakcji $R$ często nie przekracza $1\%$, oraz na problem zimnego startu dla nowych uczestników systemu. Jako rozwiązanie kompromisowe autor omawia architektury hybrydowe, które poprzez strategie łączone, sekwencyjne czy rozszerzanie cech scalają sygnały treściowe z relacjami interakcji. Ważnym wnioskiem z pracy Krawczuka jest również identyfikacja wymogów architektonicznych: wdrożenie przemysłowe na platformach VOD wymaga utrzymania niskiego czasu odpowiedzi przy jednoczesnej skalowalności. Podsumowując ewaluację, autor dowodzi, że w serwisach wideo kluczowe znaczenie ma kolejność prezentowanych propozycji na ekranie, w związku z czym metryki rankingowe Top-$K$, takie jak średnia odwrotność rangi (MRR), NDCG oraz Precision@\allowbreak k, stanowią jedyne adekwatne kryterium oceny modeli w tym środowisku.

Praca He i in. z 2017 roku skupia się na wykorzystaniu głębokich sieci neuronowych do tworzenia systemów rekomendacyjnych \footcite{heNeuralCollaborativeFilteringKwiecien032017}, wprowadzając architekturę NCF. W pierwszej części publikacji autorzy uzasadniają potrzebę zastąpienia klasycznego iloczynu skalarnego nieliniowymi warstwami neuronowymi, wykazując ograniczenia liniowych metod faktoryzacji macierzy oraz przedstawiając teoretyczne fundamenty proponowanego rozwiązania. Do przeprowadzania eksperymentów empirycznych wykorzystano zbiory danych MovieLens 1M oraz Pinterest. W przypadku MovieLens 1M oceny jawne przekształcono na dane niejawne, traktując każdą ocenę jako obecność interakcji ($y_{ui} = 1$), zbiór Pinterest poddano natomiast filtracji typu 20-core (odrzucając użytkowników posiadających mniej niż 20 wykonanych akcji). W obu zbiorach zastosowano procedurę podziału \textit{Leave-One-Out}, w której ostatnią interakcję odłożono do testów, przedostatnią do walidacji, a pozostałe do uczenia, łącząc w zestawie testowym każdy pozytywny obiekt z 99 losowymi próbkami negatywnymi. W celach porównawczych zestawiono architekturę NCF z szeregiem znanych metod referencyjnych: niepersonalizowanym modelem bazowym popularności ItemPop, modelem opartym na sąsiedztwie obiektów ItemKNN, faktoryzacją macierzy z optymalizacją błędu BPR oraz zoptymalizowaną punktowo faktoryzacją eALS (\foreignlanguage{english}{\textit{efficient Alternating Least Squares}}). Ewaluację przeprowadzono na podstawie dwóch metryk jakości rankingu Top-10: wskaźnika trafień ($HR@10$), mierzącego obecność elementu testowego wśród pierwszych 10 propozycji, oraz wskaźnika $NDCG@10$, uwzględniającego dokładną pozycję trafienia na liście. Wyniki jednoznacznie potwierdziły wyższą skuteczność hybrydowej architektury NeuMF (łączącej liniowy wariant uogólnionej faktoryzacji macierzy, ang. \foreignlanguage{english}{\textit{Generalized Matrix Factorization}} -- GMF, z nieliniową siecią MLP) w stosunku do badanych modeli odniesienia. Wykazano ponadto, że zwiększanie głębokości sieci MLP poprawia celność predykcji, a wstępne uczenie parametrów składowych przed ich fuzją pozytywnie wpływa na zbieżność i końcową skuteczność algorytmu.

Kolejny przełomowy etap w rozwoju literatury to również praca He i in., tym razem z 2020 roku \footcite{heLightGCNSimplifyingPowering2020}, wprowadzająca model LightGCN. Autorzy przeprowadzają w niej krytyczną dekonstrukcję dotychczasowych splotowych sieci grafowych (takich jak NGCF), wykazując, że przeniesienie skomplikowanych mechanizmów z klasycznych zastosowań uczenia maszynowego — takich jak nieliniowe transformacje macierzowe $\mathbf{W}$ czy nieliniowa funkcja aktywacji (np. ReLU) — jest w filtracji kolaboracyjnej zbędne i utrudnia proces uczenia. Tradycyjne sieci GCN projektowano z myślą o grafach z bogatymi opisami cech (np. sieciach cytowań z tekstami artykułów), podczas gdy w rekomendacjach obiekty i użytkownicy opisani są wyłącznie unikalnymi identyfikatorami ID. Poprzez szczegółowe eksperymenty ablacyjne autorzy udowodnili, że odrzucenie wag macierzowych oraz nieliniowości eliminuje trudności optymalizacyjne, pozwalając sieci na znacznie szybszą zbieżność i lepszą generalizację. W zamian zaproponowali architekturę LightGCN, która opiera się wyłącznie na liniowej agregacji osadzeń bezpośrednich sąsiadów na dwudzielnym grafie interakcji użytkownik-obiekt, stosując jedynie symetryczną normalizację stopni węzłów. Dodatkowo model zastępuje tradycyjne samopołączenia mechanizmem kombinacji warstw, który uśrednia reprezentacje uzyskane ze wszystkich poziomów propagacji — od zerowej warstwy identyfikatorów po wyższe poziomy połączeń wyższego rzędu. W przeprowadzonych eksperymentach na zbiorach Gowalla, Yelp2018 oraz Amazon-Book model LightGCN bezapelacyjnie pokonał dotychczasowe modele odniesienia — w tym klasyczną faktoryzację MF-BPR, głęboki perceptron NeuMF oraz bazowe sieci grafowe NGCF i GC-MC. Zastosowana uproszczona architektura pozwoliła na średni skok celności rekomendacji rankingowych o ponad 16,5\% w metryce Recall@20 oraz o 16,9\% w metryce NDCG@20 w porównaniu do pełnego modelu NGCF. Istotną obserwacją empiryczną było również zachowanie modelu przy głębszej propagacji: o ile w NGCF zwiększanie liczby warstw powyżej dwóch prowadziło do spadku celności z powodu trudności w zbieżności gradientowej, o tyle LightGCN pozwalał na stabilne wykorzystanie informacji z 3- i 4-stopniowego sąsiedztwa, stale poprawiając jakość rankingu. Ponadto całkowite odrzucenie operacji macierzowych przełożyło się na drastyczne obniżenie błędu treningowego oraz ponad dwukrotne skrócenie czasu uczenia pojedynczej epoki.

Współczesnym wyzwaniem w metodologii badań nad systemami rekomendacyjnymi stał się tzw. kryzys odtwarzalności. Znaczne rozproszenie środowisk programistycznych oraz brak standaryzowanych protokołów ewaluacji (np. stosowanie odmiennych wariantów próbkowania negatywnego czy ocenianie na ograniczonej próbie zamiast pełnego sortowania całego katalogu) powodowały, że publikowane w literaturze przewagi nowych modeli bywały trudne do obiektywnej weryfikacji. Odpowiedzią na ten problem stała się przełomowa praca Zhao i in. \footcite{zhaoRecBoleUnifiedComprehensive2020}, wprowadzająca zunifikowany szkielet badawczy RecBole. Autorzy stworzyli otwartą bibliotekę w środowisku PyTorch, która integruje ponad 70 uznanych algorytmów rekomendacyjnych (od tradycyjnych metod kolaboratywnych po zaawansowane głębokie sieci neuronowe i architektury grafowe) na podstawie spójnego formatu plików atomowych oraz jednolite pętle optymalizacyjne i ewaluacyjne. Stosując zaawansowane procedury akceleracji GPU dla predykcji rankingowych, RecBole zdefiniował nowy standard metodologiczny w dziedzinie RecSys, umożliwiając przeprowadzenie w pełni sprawiedliwego, powtarzalnego i porównywalnego benchmarku modeli wywodzących się z różnych paradygmatów technologicznych. Wpisując się w ten nurt badawczy, w niniejszej pracy wykorzystano właśnie środowisko RecBole do przeprowadzenia autorskich eksperymentów. Pozwolą one na obiektywne zestawienie opisanych modeli klasycznych (BPR), głębokich (NCF) oraz grafowych (LightGCN) na rzeczywistych danych z domeny serwisów VOD.

\subsection{Kompromisy badawcze i wyzwania wdrożeniowe w serwisach streamingowych} % REVIEWED

Przeniesienie zaawansowanych koncepcji teoretycznych, omówionych w poprzednich podrozdziałach, do komercyjnego środowiska platform VOD wiąże się z koniecznością ciągłego balansowania między czystą dokładnością matematyczną a operacyjnymi uwarunkowaniami biznesowymi i technicznymi. Projektowanie takich systemów wymaga mierzenia się z kilkoma fundamentalnymi kompromisami.

\subsubsection{Kompromis dokładności i różnorodności: wyzwanie „Długiego Ogona”}

Pierwszym kluczowym wyzwaniem jest zjawisko asymetrii popularności katalogu treści, określane w literaturze jako krzywa długiego ogona (ang. \textit{Long Tail}) \footcite{celmaMusicRecommendationDiscovery2010}. Większość interakcji użytkowników skupia się wokół wąskiej grupy najpopularniejszych hitów kinowych i serialowych, podczas gdy ogromna część katalogu platformy streamingowej generuje znikomą liczbę wyświetleń.

Klasyczne algorytmy oparte na filtrowaniu kolaboracyjnym wykazują silną tendencję do rekomendowania wyłącznie pozycji najpopularniejszych (ang. \textit{popularity bias}). Zjawisko to prowadzi do efektu „bogatych, którzy stają się jeszcze bogatsi” (\textit{rich-get-richer}) \footcite{celmaMusicRecommendationDiscovery2010}. Choć faworyzowanie powszechnie znanych hitów gwarantuje modelowi wysokie wyniki w standardowych metrykach dokładności offline, z punktu widzenia biznesowego platformy jest to wysoce nieefektywne. Tworzy to tzw. bańki filtracyjne i ogranicza widoczność niszowych produkcji, uniemożliwiając pełną monetyzację nabytej bazy filmowej \footcite{anelliTopNRecommendationAlgorithms2022}. Wybór i parametryzacja modelu muszą więc stanowić świadomy kompromis między maksymalną dokładnością rankingową a zdolnością systemu do wspierania odkrywania nowych treści.

\subsubsection{Problem zimnego startu w dynamicznych katalogach}

W przeciwieństwie do tradycyjnego rynku e-commerce ze stałym asortymentem, środowisko platform strumieniowych charakteryzuje się ciągłą rotacją katalogu. Serwisy VOD regularnie nabywają nowe licencje i publikują własne produkcje. Sytuacja ta rodzi krytyczny wdrożeniowy kompromis związany z tzw. problemem zimnego startu dla nowych obiektów.

Algorytmy kolaboracyjne, w tym modele oparte na uczeniu głębokim (NCF) oraz sieci grafowe (LightGCN), funkcjonują wyłącznie na podstawie macierzy historycznych interakcji. W momencie dodania nowego filmu nie posiada on jeszcze ocen ani historii odtworzeń. W strukturze grafu dwudzielnego staje się węzłem izolowanym (pozbawionym krawędzi), przez co model nie jest w stanie wyznaczyć dla niego odpowiedniego wektora cech (ang. \foreignlanguage{english}{\textit{embedding}}) \footcite{anelliTopNRecommendationAlgorithms2022}. Oparcie systemu wdrożeniowego wyłącznie na metodach CF skutkuje całkowitym pomijaniem nowości w generowanych rekomendacjach. W praktyce inżynieryjnej problem ten rozwiązuje się poprzez stosowanie architektur hybrydowych (uwzględniających metadane obiektów, takie jak reżyser czy gatunek) lub wprowadzanie reguł heurystycznych promujących nowe pozycje, dopóki nie zgromadzą one liczby interakcji niezbędnej do poprawnego działania modeli kolaboracyjnych.

\subsubsection{Wydajność obliczeniowa a wyzwanie czasu rzeczywistego}

Z punktu widzenia inżynierii systemów oprogramowania, zaawansowane modele matematyczne podlegają ograniczeniom wydajnościowym. Czas wygenerowania rekomendacji warunkuje czas odpowiedzi serwera i w warunkach komercyjnych nie powinien przekraczać ułamka sekundy (tzw. \textit{latency} na poziomie pojedynczych milisekund), aby zachować bezstratną płynność działania aplikacji \footcite{malinowskiAlgorytmRekomendacjiBazujacy2022}.

Złożoność strukturalna współczesnych modeli uczenia maszynowego rodzi tu poważne wyzwania wdrożeniowe:
\begin{itemize}
    \item \textbf{Modele głębokie (NeuMF/NCF):} Charakteryzują się bardzo dużą liczbą parametrów treningowych ze względu na jednoczesne uczenie gęstych osadzeń dla warstw MLP oraz GMF \footcite{heNeuralCollaborativeFilteringKwiecien032017}. Choć faza inferencji dla pojedynczego zapytania może być silnie zoptymalizowana, to dynamiczne przeliczanie nieliniowych transformacji dla milionów kombinacji bywa wysoce kosztowne w czasie rzeczywistym.
    \item \textbf{Modele grafowe (LightGCN):} Paradoksalnie, pomimo rezygnacji z nieliniowych funkcji aktywacji oraz wag, sieci GNN cechują się potężnym kosztem operacyjnym ze względu na mechanizm wieloskokowej propagacji (agregacji) informacji po całym grafie dwudzielnym \footcite{heLightGCNSimplifyingPowering2020}. Skompletowanie rekomendacji dla użytkownika wymaga pobrania cech z jego głębokiego sąsiedztwa (sieci powiązań), co stanowi potężne wąskie gardło wydajnościowe.
\end{itemize}

Rozwiązaniem tego problemu w architekturach przemysłowych jest najczęściej zastosowanie asynchronicznej dystrybucji obliczeń. Złożone algorytmy sztucznej inteligencji trenowane są okresowo w trybie wsadowym (ang. \textit{batch processing}) na gigantycznych klastrach danych, gdzie z wyprzedzeniem wyliczają statyczne wektory cech. Następnie gotowe parametry są propagowane do szybkich baz klucz-wartość w środowisku online, gdzie błyskawiczna warstwa heurystyczna odpowiada wyłącznie za ostateczne filtrowanie biznesowe (np. odrzucenie z rekomendacji filmów, na które wygasła licencja) i proste sortowanie gotowych wyników predykcji.

Opisane powyżej wyzwania, związane ze specyfiką platform VOD, problemem rzadkości macierzy ocen, zjawiskiem baniek filtracyjnych oraz narzutem obliczeniowym głębokich architektur, stanowią bezpośrednie uzasadnienie celu naukowego niniejszej pracy. Wymagają one empirycznego porównania i ewaluacji modeli o różnym stopniu skomplikowania: od klasycznego filtrowania kolaboracyjnego, poprzez architektury głębokiego uczenia, aż po grafowe sieci neuronowe. Badanie to ma na celu weryfikację skuteczności wybranych algorytmów na wysoce rozrzedzonych, rzeczywistych danych, reprezentatywnych dla współczesnej branży rozrywkowej. Cel ten został zrealizowany w kolejnych rozdziałach pracy, w których przedstawiono metodologię eksperymentu oraz analizę uzyskanych wyników.